
%(BEGIN_QUESTION)
% Copyright 2011, Tony R. Kuphaldt, released under the Creative Commons Attribution License (v 1.0)
% This means you may do almost anything with this work of mine, so long as you give me proper credit

\vbox{\hrule \hbox{\strut \vrule{} {\bf Desktop Process-øvelse} \vrule} \hrule}

En viktig funksjon i enhver prosessregulator er noe som kalles {\it utgangssporing}. Denne funksjonen gjør overgangen fra ``Auto''-modus til ``Manual''-modus enklere. Bruk en Desktop Process til å demonstrere denne funksjonen i praksis.

\vskip 10pt

En annen viktig funksjon i enhver prosessregulator er noe som kalles {\it settpunktsporing}. Denne funksjonen gjør overgangen fra ``Manual''-modus til ``Auto''-modus enklere. Bruk en Desktop Process til å demonstrere denne funksjonen i praksis.

\underbar{file i01490}
%(END_QUESTION)





%(BEGIN_ANSWER)

Når en regulator står i automatisk modus, betyr utgangssporing at den manuelle utgangsverdien følger (``sporer'') den automatiske utgangsverdien, slik at overgangen blir støtfri når regulatoren byttes til manuell modus.

\vskip 10pt

Når en regulator står i manuell modus, betyr settpunktsporing at settpunktverdien følger (``sporer'') prosessvariabelen, slik at når regulatoren byttes til automatisk modus, starter settpunktet på samme verdi som prosessvariabelen og reguleringen starter uten feil.

\vskip 10pt

Med andre ord betyr settpunktsporing at regulatoren antar at prosessen er der du ønsker at den skal være i det øyeblikket du bytter til automatisk modus.

%(END_ANSWER)





%(BEGIN_NOTES)


%INDEX% Control: setpoint tracking
%INDEX% Control: output tracking

%(END_NOTES)
